\subsection{Tangent-root attachments and the delta-five cubic-cover budget}

\begin{statement}{C.2.RAT.5}{Proved}{A locally unibranch tangent-root component gives positive Euler number}
Retain the torus-fibre, rational-normalization and specified
ruled-boundary morphism of C.2.RAT.1. Suppose $C$ is not
contracted and maps with degree $d=2$ or $3$ onto a smooth
rational curve $H$ with generic ordinary-triple transversal.
Let $q_3\in H^0(H,\operatorname{Sym}^3N^*\otimes M)$ be
its cubic normal tensor, where $N=N_{H/Y}$ and
$M=\OO_Y(S)|_H$. Suppose $q_3$ is nonzero as a tensor at
every point of $H$. Its tangent-root divisor
$R_3\subset\PP(N)$ is then finite of degree three over $H$.
Let $R_d$ be the component corresponding to $C$ over the
generic point. If $R_d$ is locally unibranch at every point,
then
\[
                           e(S)\ge d-1>0.
\tag{C.RD.1}
\]
In particular this applies when $d=2$ and the only singularities
of $R_2$ are an $A_4$ cusp or two $A_2$ cusps.
\end{statement}
\begin{proof}
The nonzero cubic has three distinct roots over the generic
point and is never identically zero on a projective fibre.
It therefore defines a reduced finite divisor $R_3$.
The function-field algebra of its roots agrees with that of
the normalized conductor curves dominating $H$: a branch
of an ordinary triple is distinguished by its tangent line.
Taking the integral closure of $\OO_H$ in that algebra gives
the exact finite-cover identification
\[
                            C\simeq R_d^\nu
                            \quad\text{over }H.
\tag{C.RD.2}
\]
Thus two distinct points of $C$ over one point $p\in H$
give distinct root points of $R_d$, because normalization
of a locally unibranch reduced curve is bijective on its
underlying points.

We check what a simple root implies for the actual
hypersurface, including its conductor. In coordinates
$H=(x,y)=0$, with parameter $z$ on $H$, the local equation
$f_0$ of $S$ belongs to $(x,y)^3$. The lower normal jets
vanish on the reduced smooth curve, and the graded pieces
of its regular-sequence ideal are locally free, which
justifies this membership. Choose the normal coordinates
so a specified simple tangent root is finite in $x/y$.
Then
\[
 f_0(y\lambda,y,z)/y^3
       =q_3(\lambda,1,z)+y\,A(\lambda,y,z).
\tag{C.RD.3}
\]
The simple root gives a holomorphic solution
$\lambda=\lambda(y,z)$ by the implicit-function theorem.
Division by the monic linear factor yields an actual
factorization
\[
 f_0=(x-y\lambda(y,z))\,G(x,y,z),\qquad
 G|_{x=y\lambda(y,z)}=y^2 u(y,z),\quad u(0)\ne0.
\tag{C.RD.4}
\]
The unit is the derivative of the cubic with respect to
its root coordinate, at the chosen simple root. Hence
this normalization branch is smooth and its conductor
in $S$ is exactly twice the curve lying over $H$.
There is no vertical conductor component there, and
the minimal resolution $\pi$ is an isomorphism nearby.

Two distinct $C$-points over $p$ cannot both receive a
vertical boundary attachment. Their root points are
distinct by (C.RD.2). If both could receive attachments,
(C.RD.4) would make both roots multiple, requiring total
cubic degree at least four. Thus each point of $H$
receives at most one attachment from $C$.

There are at most $3-d$ generic normalization branches
left to be supplied by vertical strict conductor
components. In the notation of C.2.RAT.3 this gives
$r_{\rm dom}\le3-d$. At least
$K\ge r-N_0-r_{\rm dom}$ original positive vertical
boundaries therefore survive with an image different
from $H$. Their $K$ points of $C$ give $K$ distinct
normalization-incidence edges at $H$, by the preceding
paragraph. Equation (C.RA.24) gives $n_G\ge K$, and
the exact Euler decomposition gives
\[
 e(S)\ge2-r+N_0+K
       \ge2-r_{\rm dom}\ge d-1.
\]
This proves the assertion. The final examples are locally
unibranch by their displayed cusp equations.
\end{proof}

\begin{statement}{C.2.RAT.6}{Known, with exact derivation}{Cubic suspensions of plane curves with delta at most five}
Let $r(x,y)=0$ be a reduced plane-curve germ with
$\delta(r)\le5$, and let $X_r$ be the isolated surface
singularity $z^3=r(x,y)$. Then $p_g(X_r)\le1$ and the
first rational Betti number of its link is at most two.
The latter is two precisely for plane type $A_5$ or $D_4$.
For a collection of these plane germs with total delta
invariant at most five,
\[
             \sum p_g(X_r)\le1,
           \qquad \sum b_1(\operatorname{link}X_r)\le2.
\tag{C.RD.5}
\]
\end{statement}
\begin{proof}
We first check that the finite normal-form list is exhaustive.
The delta formula under a point blowup shows that a plane
germ of multiplicity at least four has delta at least six.
For an irreducible branch of multiplicity two,
Weierstrass preparation, completion of the square and a
root of a unit give $y^2-x^{2k+1}$, of delta $k\le5$.
For an irreducible branch of multiplicity three, the
Puiseux conductor formula gives delta at least six
unless its first characteristic exponent is four or
five. A term with exponent divisible by three is removed
by subtracting a power of the first coordinate. The
remaining two possibilities have relatively prime
exponents $(3,4)$ or $(3,5)$, and no later characteristic
pair. Their analytic normal forms are $E_6$ and $E_8$.
This last analytic comparison is also the zero-parameter
case of the semiquasihomogeneous normal-form theorem:
the bases of the two gradient algebras are respectively
$x^i y^j$ with $i\le2,j\le1$, and with
$i\le3,j\le1$, after interchanging variables if necessary.
Every basis monomial has weight less than one, so there
is no higher-weight term left in the normal form; see
\cite[\S1, I.1.0, and \S2]{Arnold1976NormalForms}.

For a reducible germ the exact normalization sequence
for a union gives
\[
 \delta\Bigl(\bigcup_i D_i\Bigr)
       =\sum_i\delta(D_i)+\sum_{i<j}I(D_i,D_j).
\tag{C.RD.6}
\]
Four branches already contribute at least six.
With three branches, a singular branch would contribute
at least $1+2+2+1=6$. Thus all three are smooth.
They are either pairwise transverse, or exactly one
pair is tangent with contact two or three: if all
three had a common tangent, the sum would again be
at least six. Straightening the transverse branch,
subtracting the mean of the other two graph functions
and taking a root of their unit difference gives
$D_4,D_6,D_8$ respectively.

With two smooth branches, subtracting their mean graph
and rescaling the transverse coordinate gives
$A_{2k-1}$ with contact $k\le5$. With two singular
branches the two delta contributions and their
intersection multiplicity give at least $1+1+4=6$.
The remaining possibility is one smooth branch and
one cusp. A multiplicity-three cusp would already
give $3+3=6$. A multiplicity-two cusp of delta $k$
has parameterization $(t^2,t^{2k+1})$.
Its intersection with a smooth branch is either two,
or, for the common tangent, one of
$4,6,\ldots,2k,2k+1$.
The inequality $k+I\le5$ leaves a transverse smooth
branch with $k=1,2,3$, giving $D_5,D_7,D_9$, or the
ordinary cusp with its tangent smooth branch, giving
$E_7$. The actual analytic forms are again those
obtained by the same graph changes and square
completion, or the zero-parameter normal forms in the
cited theorem. The resulting complete list is
\[
              A_1,\ldots,A_{10},\qquad
              D_4,\ldots,D_9,\qquad E_6,E_7,E_8.
\tag{C.RD.7}
\]
An analytic equivalence of the plane equation up to
a unit gives an analytic equivalence of the cubic
suspensions: take a local cube root of that unit and
rescale $z$. Thus these are the precise surface germs
whose invariants must be computed.

The normal forms, with weights normalized so the
equation has degree one, are
\[
\begin{array}{c|c|c|c}
 \text{plane type}&r(x,y)&(w_x,w_y)&\delta(r)\\\hline
 A_n&x^2+y^{n+1}&(1/2,1/(n+1))&\lfloor(n+1)/2\rfloor\\
 D_n&x^2y+y^{n-1}&((n-2)/(2(n-1)),1/(n-1))
                         &\lfloor(n+2)/2\rfloor\\
 E_6&x^3+y^4&(1/3,1/4)&3\\
 E_7&x^3+xy^3&(1/3,2/9)&4\\
 E_8&x^3+y^5&(1/3,1/5)&4.
\end{array}
\tag{C.RD.8}
\]
The third weight of the suspension is $1/3$.
The monomial-spectrum formula and its geometric-genus
interval, or the direct adjunction-ideal calculation
\cite[Theorem 2.1.1 and Corollary 2.1.3]{MerleTeissier1977},
give
\[
 p_g(X_r)=\#\{(i,j,k)\in\mathbb Z_{\ge0}^3:
          (i+1)w_x+(j+1)w_y+(k+1)/3\le1\}.
\tag{C.RD.9}
\]
These particular forms are Newton nondegenerate;
this follows directly by checking that their face
polynomials have no common zero of their derivatives
in the relevant algebraic torus. In (C.RD.9) either
there is no monomial, for $A_1,\ldots,A_4$, or the
constant monomial alone satisfies the inequality.
In particular a nonzero geometric genus requires
delta at least three.

For clarity, the link calculation uses the actual
Milnor-fibre monodromy rather than a homology-sphere
assumption. For a monomial $x^i y^j$ in the plane
Jacobian algebra, its monodromy exponent modulo
integers is $(i+1)w_x+(j+1)w_y$. The two reduced
eigenvalues of $z^3$ are the two primitive cube
roots. The Thom--Sebastiani tensor comparison
therefore makes the eigenvalue-one dimension for
the suspension equal to the number of plane
exponents congruent to $1/3$ or $2/3$.
The isolated-surface Milnor fibre is simply
connected. Its pair sequence and Poincare--Lefschetz
duality identify the rational first homology of
its boundary with the radical of the middle
intersection pairing. The variation comparison
identifies that radical dimension with the
eigenvalue-one dimension. These weighted homogeneous
monodromies are semisimple; the monomial spectrum
and suspension formula are given in
\cite[pp. 165--167]{Steenbrink1989Spectrum}.

For $A_n$ the plane Jacobian basis is
$1,y,\ldots,y^{n-1}$; for $D_n$ it is
$1,y,\ldots,y^{n-2},x$; for $E_6,E_8$ it is the
rectangular basis from the preceding paragraph.
For $E_7$ it is $1,y,\ldots,y^4,x,xy$.
Direct substitution gives
\[
\begin{array}{c|c|c}
 \text{type}&p_g(z^3-r)&b_1(\operatorname{link}(z^3-r))\\\hline
 A_1,\ldots,A_4&0&0\\
 A_5&1&2\\
 A_6,\ldots,A_{10}&1&0\\
 D_4&1&2\\
 D_5,\ldots,D_9&1&0\\
 E_6,E_7,E_8&1&0.
\end{array}
\tag{C.RD.10}
\]
All weights, the individual newly required
$A_9,A_{10},D_8,D_9$ rows and every unordered
collection with total delta at most five are
checked with exact rational arithmetic in
\path{code/A_cubic_delta5.py} and its saved output.
Since each nonzero entry in either invariant
costs delta at least three, at most one such
entry occurs in a collection of total delta
at most five. This proves (C.RD.5).
\end{proof}

\begin{statement}{C.2.RAT.7}{Proved}{Every primitive rational double ordinary-triple cover has nonnegative Euler number}
Retain the torus-fibre, rational-normalization and specified
ruled-boundary morphism of C.2.RAT.1. Suppose $C$ is not
contracted, its image $H$ has generic ordinary-triple
transversal, and $C\to H^\nu$ has degree two. Suppose
the cubic normal tensor along $H$ has no common zero of
its coefficients. Then $e(S)\ge0$. This holds for every
splitting type of the normal bundle $N_{H/Y}$.
\end{statement}
\begin{proof}
Statement A.5.19 makes $H$ smooth: over a singular
rational ordinary-triple image every connected normalized
conductor cover has degree one. Thus $H\simeq\PP^1$.
Set $N=N_{H/Y}$, $W=\PP(N)$, and
$\xi=c_1(\OO_W(1))$, where $\PP(N)$ parameterizes lines.
The ambient Chern-class hypothesis gives
$\deg N=-2$ and $\OO_Y(S)|_H\simeq\OO_H$.
The cubic is an actual section of $\OO_W(3)$.
The assumed absence of a common coefficient zero makes
its root divisor finite of degree three over $H$:
\[
 R=R_1+R_2\sim3\xi,\qquad
 \xi^2=2,\qquad K_W=-2\xi,\qquad p_a(R)=4.
\tag{C.RD.11}
\]
Here $R_1$ is a section and hence is smooth rational;
$R_2$ has normalization $C\simeq\PP^1$, by the exact
integral-closure comparison (C.RD.2). Both are reduced
and they are different. The normalization sequence gives
\[
                  \sum_p\delta_p(R)=5.
\tag{C.RD.12}
\]
This equality also holds if the two components are
disjoint: the arithmetic-genus convention is
$p_a(R)=1-\chi(\OO_R)$ and the normalization has Euler
characteristic two.

Form the cyclic triple cover
$p_0:X_0\to W$ with algebra
$\OO_W\oplus\OO_W(-1)\oplus\OO_W(-2)$ and multiplication
given by the cubic. It is normal because its branch
divisor $R$ is reduced, and is connected because a
meridian around a smooth branch point permutes its
three sheets cyclically. Resolve its isolated
singularities and its reduced boundary to obtain
$X$, with $D=X\setminus U$, where
$U=\{v\in N:q_3(v)=1\}$ is the smooth normal-cone
level surface. The isomorphism over $W\setminus R$
is the exact radial-vector/dual-tautological-coordinate
comparison (C.D4.6).

The projective-bundle cohomology calculation gives
\[
 H^1(W,\OO_W)=H^1(W,\OO_W(-1))
             =H^1(W,\OO_W(-2))=0.
\tag{C.RD.13}
\]
Indeed both relative cohomologies of $\OO_W(-1)$
vanish, while $p_*\OO_W(-2)=0$ and
$R^1p_*\OO_W(-2)=\det N=\OO_{\PP^1}(-2)$.
Thus $H^1(X_0,\OO_{X_0})=0$ for every normal-bundle
splitting. Resolution Leray and C.2.RAT.6 give
$q(X)\le\sum p_g(X_0,x)\le1$, so $b_1(X)\le2$.

We spell out the boundary calculation with its number
of connected components retained. Write
\[
 \nu=\sum_p(r_p(R)-1),\qquad
 c=b_0(R),\qquad
 \lambda=\sum_x b_1(\operatorname{link}(X_0,x))\le2.
\]
Normalization of $R$ gives $e(R)=4-\nu$.
Since $U\to W\setminus R$ is etale of degree three,
$e(U)=3\nu$. The resolved boundary has the same
number $c$ of connected components. Replacing each
singular point of $R$ by the connected exceptional
divisor of its cubic suspension gives, by the
Mayer--Vietoris sequence with its individual branch
attachment points,
\[
                 b_1(D)=\nu-2+c+\lambda.
\tag{C.RD.14}
\]
The exceptional first Betti number is the link's
first Betti number: the exceptional intersection
matrix in its plumbed neighbourhood is invertible
over $\Q$. This is precisely the local comparison
proved with (C.D4.7). Additional boundary-resolution
blowups add rational trees and leave (C.RD.14)
unchanged.

Let $d_D$ be the number of irreducible components
of $D$, and let $V\subset H_2(X,\Q)$ be the span
of their classes. The exceptional classes are
independent. Pushing a relation to $W$ and pulling
back a relation from $R_1,R_2$ shows that
\[
 \varepsilon:=d_D-\dim V
      =2-\dim\operatorname{span}\{[R_1],[R_2]\}
      \in\{0,1\}.
\]
This proves the relation count rather than
assuming the two branch classes are independent.
The space $V$ contains the pullback of $R$,
whose square is $3R^2=54>0$. The Hodge index
theorem therefore makes its intersection pairing
nondegenerate. The divisor Gysin sequence gives
$b_1(U)=b_1(X)+\varepsilon\le3$; its degree-two
part identifies the image of $H_2(U)$ in
$H_2(X)$ with $V^\perp$. Consequently the rank
of the compact intersection pairing on $U$ is
$b_2(X)-d_D+\varepsilon$.
Now use $e(X)=2-2b_1(X)+b_2(X)$ and
$e(D)=c-b_1(D)+d_D$. Equations (C.RD.14) and
$e(U)=3\nu$ yield the exact identity
\[
\begin{split}
 \operatorname{rank}I_U
   &=e(U)-2+2b_1(X)+c-b_1(D)+\varepsilon\\
   &=2\nu+2b_1(X)-\lambda+\varepsilon
     \ge2\nu-2.
\end{split}
\tag{C.RD.15}
\]
In particular this formula includes disjoint root
components; no connected-boundary assumption has
been introduced.

If $\nu\ge2$, (C.RD.15) gives a nonzero compact
intersection pairing on a surface with $b_1(U)\le3$.
The compact-core transport (C.D4.2) embeds that
pairing into an actual nearby torus. Its induced
first-homology image has dimension at most three,
so its second-homology image has zero wedge pairing,
as established immediately after (C.D4.2). This
contradiction proves $\nu\le1$.

Finally suppose a value of $H$ receives vertical
boundary attachments at both its $C$-preimages.
If their tangent-root points in $R_2$ were distinct,
at least one would be a simple root of the cubic,
and (C.RD.3)--(C.RD.4) would give pure conductor
there, a contradiction. They therefore map to a
single multibranch point of $R_2$. Distinct values
give distinct such points, each contributing at
least one to $\nu$. Hence there is at most one
value with two attachments. With
$K\ge r-N_0-r_{\rm dom}$ surviving original
vertical boundaries, at least $K-1$ distinct
normalization-incidence edges remain at $H$.
There is only one remaining generic triple branch,
so $r_{\rm dom}\le1$. The graph comparison
(C.RA.24) and exact Euler decomposition now give
\[
 e(S)\ge2-r+N_0+(K-1)
                         \ge1-r_{\rm dom}\ge0.
\tag{C.RD.16}
\]
This completes the proof.
\end{proof}

\begin{statement}{C.2.RAT.8}{Proved}{A single common cubic zero in the rational double-cover case}
Retain the torus-fibre, rational-normalization and specified
ruled-boundary morphism of C.2.RAT.1. Suppose $C$ is not
contracted, its image has generic ordinary-triple transversal,
and $C\to H^\nu$ has degree two. If the common zero divisor
of the cubic normal coefficients has degree one, then
$e(S)\ge0$.
\end{statement}
\begin{proof}
As in C.2.RAT.7, $H$ is smooth rational. On $W=\PP(N)$
write the complete root divisor as
\[
 B=R_1+R_2+F\sim3\xi,
 \qquad B_h=R_1+R_2\sim3\xi-F.
\tag{C.RD.17}
\]
Here $F$ is the ruling fibre over the unique simple
common zero. The horizontal divisor $B_h$ is finite
over $H$, $R_1$ is a section, and the normalization of
$R_2$ is the rational degree-two cover $C$.
Since $K_W=-2\xi$, $\xi^2=2$, and $\xi F=1$,
\[
 p_a(B_h)=2,
 \qquad \sum_p\delta_p(B_h)=3,
 \qquad I(F,B_h)=3.
\tag{C.RD.18}
\]
The full divisor $B$ has total delta six. We will
use its particular decomposition (C.RD.17), not a
statement about arbitrary delta-six plane germs.

\paragraph{The local cubic-suspension genus along $F$.}
At a point of $F$, put $i=I_p(F,B_h)$. If $i=1$,
the germ of $B$ is an ordinary node. If $i=2$,
the horizontal germ is either smooth with contact
two with $F$, two smooth branches transverse to
$F$, or a multiplicity-two cusp transverse to $F$.
The bound $\delta_p(B_h)\le3$ makes the resulting
plane types respectively
\[
 A_3,\qquad D_{2k+2}\ (1\le k\le3),\qquad
                 D_{2k+3}\ (1\le k\le3).
\tag{C.RD.19}
\]
For two smooth branches $k$ is their contact;
for the cusp $k$ is its delta invariant.
The exact analytic identifications are the graph
straightening and square-completion comparisons
in C.2.RAT.6. Their cubic suspensions have
geometric genus zero or one by (C.RD.10).

If $i=3$, the point contains $R_1$ and all the
intersection of $R_2$ with $F$. The section $R_1$
is transverse to $F$. A smooth branch of $R_2$
with intersection two is tangent to $F$ and
transverse to $R_1$, giving $D_6$. If $R_2$ has
two local branches, both are smooth and transverse
to $F$. The horizontal delta bound three forces
all three horizontal branches to be pairwise
transverse. The full germ is therefore four
distinct transverse smooth branches.
The last possibility is one singular branch
of $R_2$. Its multiplicity is two, since its
intersection with $F$ is two. Its intersection
with $R_1$ is at least two; hence (C.RD.18)
forces it to be an ordinary cusp, transverse
to both smooth branches $F$ and $R_1$.
These alternatives exhaust $i=3$.

For four transverse branches, the degree-four
initial plane polynomial has four distinct
linear factors. Its cubic suspension is Newton
nondegenerate with weights $(1/4,1/4,1/3)$;
the only positive lattice point on or below
its compact Newton face is $(1,1,1)$. Its
geometric genus is therefore one by
\cite[Theorem 2.1.1 and Corollary 2.1.3]{MerleTeissier1977}.

We give the precise same calculation for the cusp
and two transverse branches, retaining higher
terms of the actual equation. Straighten $F,R_1$
to coordinate axes and make a linear change so
the cusp tangent is $u=0$. Up to units, the
equation is then
\[
 r(u,v)=v(u-v)\bigl(u^2+c v^3+
           u\,O((u,v)^2)+O((u,v)^4)\bigr),
                  \qquad c\ne0.
\tag{C.RD.20}
\]
The nonzero radical cubic is the ordinary-cusp
condition. The compact Newton facets of
$z^3-r(u,v)$ have normalized inward vectors
\[
             (1/4,1/4,1/3),\qquad
             (3/10,1/5,1/3).
\tag{C.RD.21}
\]
Their plane edges have vertex sets
$\{(3,1),(2,2)\}$ and
$\{(2,2),(0,5)\}$. The higher terms in
(C.RD.20) are strictly above these faces.
The first edge polynomial is a nonzero multiple
of $u^2v(u-v)$; the second is a nonzero multiple
of $v^2(u^2+c v^3)$. Both have smooth zero sets
in $(\C^*)^2$. Each two-dimensional compact
face additionally contains the term $z^3$,
whose $z$-derivative is nonzero on $(\C^*)^3$.
This proves the required Newton nondegeneracy.
Every noncompact face inequality is strictly
satisfied by positive integral points; the only
additional one beyond coordinate positivity is
$v+z/3\ge1$.
For either vector in (C.RD.21), the point
$(1,1,1)$ has weight $5/6$, and increasing any
one coordinate makes its weight greater than
one. Thus it is the only positive lattice
point outside the interior of the Newton
polyhedron. The same adjunction theorem gives
$p_g(z^3-r)=1$.

This proves that every suspension singularity
over $F$ has geometric genus at most one.
A nonzero value requires $i\ge2$ and
$\delta_p(B_h)\ge1$: the cases with horizontal
delta zero were the ordinary node and $A_3$.
Since the total intersection with $F$ is three,
there is at most one such nonzero value over
$F$. Away from $F$, C.2.RAT.6 shows that a
nonzero geometric genus costs horizontal delta
at least three. It cannot coexist with a
nonzero value over $F$, by (C.RD.18).
Consequently the cyclic cover $X_0$ branched
over $B$ satisfies
\[
                        \sum_x p_g(X_0,x)\le1.
\tag{C.RD.22}
\]

We also obtain the required link bound without
assuming a delta-six monodromy table. For a
normal surface singularity, let $E$ be the
reduced exceptional divisor of a good resolution.
The restriction maps from infinitesimal
neighbourhoods of $E$ onto $E$ are surjective
on $H^1$, since the kernels have one-dimensional
support and no $H^2$. Formal functions therefore
gives $h^1(E,\OO_E)\le p_g$.
For an SNC exceptional curve,
$b_1(E)=2\sum g_i+\beta\le
2(\sum g_i+\beta)=2h^1(E,\OO_E)$.
The invertible exceptional intersection matrix
identifies this $b_1(E)$ with the link's first
Betti number, as in the proof of C.2.RAT.6.
Equation (C.RD.22) yields
\[
              \lambda:=\sum_x b_1(\operatorname{link}(X_0,x))
                            \le2.
\tag{C.RD.23}
\]

\paragraph{The compact pairing and the attachment count.}
The cyclic algebra and projective-bundle cohomology
are unchanged from (C.RD.13), so its projective
resolution $X$ has $b_1(X)\le2$.
The divisor $B$ is connected, since $F$ meets
both horizontal components. Its three classes
span a two-dimensional space: the section
$R_1$ and the fibre $F$ are independent.
Thus the boundary-relation number is exactly
$\varepsilon=1$, and $b_1(U)=b_1(X)+1\le3$.
Put $\nu=\sum_p(r_p(B)-1)$.
Normalization of its three rational components
and the resolved-boundary comparison give
\[
 e(U)=3(\nu-2),\qquad
 b_1(D)=\nu-2+\lambda,
 \qquad
 \operatorname{rank}I_U
       =2\nu-4+2b_1(X)-\lambda\ge2\nu-6.
\tag{C.RD.24}
\]
The rank identity is the same Gysin, Hodge-index
and Euler calculation as (C.RD.15), with all
three boundary components and their one relation
included. Compact-core transport into the
nearby torus therefore forces $\nu\le3$.

Let $p\in H$ be the common zero value and
$a=\#(C\to H)^{-1}(p)\in\{1,2\}$.
At every point of $F\cap B_h$, the full divisor
has the branch of $F$ together with the
individual horizontal normalization branches.
It follows exactly that
\[
        \nu_F:=\sum_{x\in F}(r_x(B)-1)=1+a.
\tag{C.RD.25}
\]
The contribution one is the section $R_1$;
the other $a$ branches are the normalization
points of $R_2$ above $p$. This identity retains
coinciding tangent roots and ramification.

The loss of distinct horizontal anchor values
at $p$ is at most $a-1=\nu_F-2$. Away from
$p$, the cubic is primitive. The argument at
the end of C.2.RAT.7 makes each duplicate
attachment there a multibranch point of $R_2$,
costing at least one in
$\nu_{\rm out}=\nu-\nu_F$. Hence the total
loss $\ell$ of distinct normalization-incidence
edges satisfies
\[
                    \ell\le\nu_F-2+\nu_{\rm out}
                              =\nu-2\le1.
\tag{C.RD.26}
\]
Finally $r_{\rm dom}\le1$ and
$K\ge r-N_0-r_{\rm dom}$. Projecting these
$K$ surviving attachments to $H$ retains at
least $K-\ell$ distinct edges. Equation
(C.RA.24) and the Euler decomposition give
\[
                 e(S)\ge2-r_{\rm dom}-\ell\ge0.
\tag{C.RD.27}
\]
The two Newton facets, their exact lower
lattice-point counts and the attachment
budgets are recorded in
\path{code/A_cubic_scalar_one.py} and its output.
\end{proof}

\begin{statement}{C.2.RAT.9}{Proved}{Every rational double ordinary-triple cover has nonnegative Euler number}
Retain the torus-fibre, rational-normalization and specified
ruled-boundary morphism of C.2.RAT.1. Suppose $C$ is not
contracted, its image has generic ordinary-triple transversal,
and $C\to H^\nu$ has degree two. Then $e(S)\ge0$,
including all common zeros and their multiplicities in
the cubic normal tensor.
\end{statement}
\begin{proof}
Again A.5.19 makes $H$ smooth rational. Let $A$ be the
common zero divisor of the cubic coefficients and put
$v=\deg A$. The cases $v=0,1$ are C.2.RAT.7--8.
We first verify that $v\le2$. Write
$N=\OO(a)\oplus\OO(-2-a)$ with $a\ge-1$.
For $a\ge0$ and $v>0$, the coefficient of the cube
of the first normal coordinate has degree $-3a-v$
and vanishes. Thus the other coordinate is a factor.
The remaining quadratic is irreducible over the
function field, because it describes the connected
degree-two cover. Its first-coordinate-square
coefficient must be nonzero and has degree
$2-a-v$. This gives $a+v\le2$.
For $a=-1$, on $\PP^1\times\PP^1$ the section
$R_1$ and quadratic component have bidegrees
$(1,d)$ and $(2,3-v-d)$, with $d\ge0$ and the
first entry denoting degree over $H$. An integral
quadratic component requires $3-v-d\ge1$:
a divisor of bidegree $(2,0)$ is a union of
constant sections over $\C$. Thus $v\le2$ here
as well.

It remains to consider $v=2$. With the notation
of C.2.RAT.7, the horizontal divisor is
\[
 B_h=R_1+R_2\sim3\xi-2F,
       \qquad p_a(B_h)=0,
       \qquad \sum_p\delta_p(B_h)=1.
\tag{C.RD.28}
\]
Both normalization components are rational.
Consequently its only possible singularity is
one ordinary node or one ordinary cusp. An
intersection between $R_1$ and $R_2$, if present,
is the node. Otherwise a node or cusp belongs
to $R_2$. Every other point of the horizontal
divisor is smooth.

\paragraph{Two distinct common zero values.}
First let $A=p+q$ with $p\ne q$. The full branch
divisor is reduced:
$B=R_1+R_2+F_p+F_q\sim3\xi$.
Use the ordinary cyclic algebra of (C.RD.13).
The local analysis in the proof of C.2.RAT.8
applies at each of the two fibres. A nonzero
local geometric genus there requires positive
horizontal delta. The total horizontal delta
in (C.RD.28) is one, and the two fibres are
disjoint. Away from these fibres the local
plane delta is at most one and its cubic
suspension is rational. Therefore
\[
 P:=\sum_x p_g(X_0,x)\le1,\qquad
 \lambda:=\sum_x b_1(\operatorname{link}(X_0,x))\le2.
\tag{C.RD.29}
\]

In this case the irregularity is zero, not
merely at most one. The finite-cover dualizing
formula gives
$\omega_{X_0}=p_0^*(\omega_W\otimes\OO_W(2))
\simeq\OO_{X_0}$, and (C.RD.13) gives
$H^1(X_0,\OO)=0$. Thus $\chi(\OO_{X_0})=2$.
For a resolution $\rho:X\to X_0$, local duality
and Grauert--Riemenschneider vanishing identify
$\rho_*\omega_X\subset\omega_{X_0}$ with an
ideal of colength $p_g(X_0,x)$ at each singular
point. Explicitly this is the dual of the
finite-length $R^1\rho_*\OO_X$ in the surface
duality sequence. If $P>0$, the global generator
of $\omega_{X_0}$ does not belong to that ideal
at a point with positive colength. No nonzero
global canonical form on $X$ can then exist,
because its pushforward would be a nonzero
constant multiple of that generator.
Resolution Leray now gives
\[
 \chi(\OO_X)=2-P,
 \qquad q(X)=P-1\quad(P>0),
 \qquad q(X)=0\quad(P=0).
\tag{C.RD.30}
\]
Since $P\le1$, this proves $b_1(X)=0$.

The four branch classes span the two-dimensional
space generated by a section and a fibre, so
$\varepsilon=2$ and $b_1(U)=2$.
The branch divisor and resolved boundary are
connected. For
$\nu=\sum_x(r_x(B)-1)$, normalization and the
exact pairing calculation give
\[
 e(U)=3(\nu-4),\qquad b_1(D)=\nu-3+\lambda,
 \qquad \operatorname{rank}I_U=2\nu-8-\lambda.
\tag{C.RD.31}
\]
The compact torus comparison forces this rank
to be zero; otherwise its nondegenerate compact
pairing would embed into a rank-at-most-two
first-homology image in the torus. Hence
$\nu=4+\lambda/2\le5$.

Let $a_p,a_q\in\{1,2\}$ be the numbers of
$C$-preimages of these two common values.
The exact branch count (C.RD.25) gives
$\nu_{F_p}+\nu_{F_q}=2+a_p+a_q$.
The loss of distinct attachment values on
these fibres is at most $a_p+a_q-2$.
Away from them each loss costs a multibranch
point of $R_2$. Consequently
\[
                    \ell\le\nu-4\le1.
\tag{C.RD.32}
\]
The graph count gives
$e(S)\ge2-r_{\rm dom}-\ell\ge0$.

\paragraph{One double common zero value.}
Now let $A=2p$. Normalize the cyclic cover
with branch multiplicities $R_1+R_2+2F_p$.
Its exact character algebra is
\[
 (p_0)_*\OO_{X_0}
       =\OO_W\oplus L^{-1}\oplus L^{-2}(F_p),
                  \qquad L=\OO_W(1).
\tag{C.RD.33}
\]
The additional integral generator at $F_p$
is $w^2/s$ in $w^3=s^2u$. Its divisorwise
character line extends reflexively over the
smooth base, which proves the sheaf identity
at the finitely many intersections as well;
this is also (C.DD.16). Relative duality gives
$R^1p_*L^{-2}(F_p)=\OO_H(-1)$.
Thus $H^1(X_0,\OO_{X_0})=0$.

We check that every singularity of this
normalized cover is rational, including the
horizontal node and cusp allowed by (C.RD.28).
For this purpose write a cyclic branch divisor
as $D_1+2D_2$, with root line $L$.
At a point blowup $\sigma$ let
$m_i=\operatorname{mult}D_i$,
$n=m_1+2m_2=3k+j$, $0\le j<3$, and let $E$
be the new exceptional curve. Keeping the
same function field gives root line
$L'=\sigma^*L(-kE)$ and new exceptional
branch weight $j$. Its two character sheaves
are exactly
\[
\begin{split}
 (L')^{-1}&=\sigma^*(L^{-1})(kE),\\
 (L')^{-2}(D'_2)
   &=\sigma^*(L^{-2}(D_2))
       \bigl((2k-m_2+\mathbf1_{j=2})E\bigr).
\end{split}
\tag{C.RD.34}
\]
Here $D'_2=\sigma^*D_2-m_2E+\mathbf1_{j=2}E$.
These identities follow by replacing the
cyclic root by its quotient by the $k$th
power of a local exceptional equation.

At a smooth horizontal branch the contact
with $F_p$ is one or two. The weight-$(2,1)$
contact calculation (C.DD.18)--(C.DD.20)
separates the branches with character
increments $(1,1)$ at each blowup.
If the horizontal node lies on $F_p$, there
are three smooth local branches of weights
$(2,1,1)$. They are transverse except possibly
for contact two between $F_p$ and the
quadratic branch in the case of an intersection
$R_1\cap R_2$. A blowup has $n=4,m_2=1$,
so both character increments in (C.RD.34)
are one. In the transverse case it separates
the three branches. In the contact-two case
one remaining ordinary triple has the same
weights $(2,1,1)$; the same blowup separates it.

If the horizontal cusp lies on $F_p$, its
local form is a weight-one cusp transverse
to the weight-two fibre:
\[
                        w^3=s^2(x^2-s^3).
\tag{C.RD.35}
\]
The first blowup has $n=4,m_2=1$ and
increments $(1,1)$. The strict cusp becomes
smooth with contact two with the new
weight-one exceptional curve. Blow up that
weight-$(1,1)$ tangency; here $n=2,m_2=0$
and the increments are $(0,1)$.
The remaining ordinary triple has weights
$(1,1,2)$, so a final blowup has increments
$(1,1)$. The resulting branch has normal
crossings. At a crossing of weights $a,b$
in $\{1,2\}$, normalization of
$w^3=x^a y^b$ is the quotient of $\C^2$
under the subgroup
\[
 \{(\zeta^i,\zeta^j):ai+bj=0\pmod3\},
 \qquad x=u^3,\quad y=v^3,\quad w=u^av^b.
\]
These are the rational cyclic quotient
singularities $1/3(1,1)$ or $1/3(1,2)$.
Their standard resolutions have respectively
one $(-3)$-curve or a two-curve $(-2)$ chain,
and $R^1\OO=0$. The horizontal node or cusp
away from $F_p$ likewise has a rational
cubic suspension by (C.RD.10).

For completeness, none of the preceding
base blowups adds a geometric-genus defect.
Every exceptional character increment in
(C.RD.34) was zero or one. A point blowup
has $\sigma_*\OO(E)=\OO$ and
$R^1\sigma_*\OO(E)=0$, from
$\OO(E)|_E=\OO_{\PP^1}(-1)$ and the
structure-sheaf direct images; the same
holds for increment zero. Projection formula
therefore makes the first direct image of
each normalized character algebra vanish.
The normalized cyclic covers over successive
base blowups map properly and birationally
to the preceding cover through their same
function field. Finite direct image and
Leray identify their $R^1\OO$ with the
just computed eigensheaf direct images;
faithfulness of finite pushforward gives
$R^1\OO=0$. Resolving the final rational
cyclic quotients preserves this property.
This proves the claimed local rationality.

We consequently have $b_1(X)=0$ and
$\lambda=0$. The reduced branch support
$B=R_1+R_2+F_p$ has three rational
normalization components, is connected,
and has class rank two. Thus
\[
 b_1(U)=1,\qquad e(U)=3(\nu-2),\qquad
 b_1(D)=\nu-2,\qquad
          \operatorname{rank}I_U=2\nu-4.
\tag{C.RD.36}
\]
Again compact transport forces rank zero,
hence $\nu=2$.
Equation (C.RD.25) then gives exactly one
$C$-preimage over the common value and no
multibranch point of $R_2$ away from it.
There is no loss of distinct attachment
values. The graph inequality yields
\[
                    e(S)\ge2-r_{\rm dom}\ge1.
\tag{C.RD.37}
\]
This proves the degree-two-common-divisor
case and, together with C.2.RAT.7--8 and
the degree bound at the start, every case
in the statement.
The additional weighted node/cusp charts,
their exact character increments and both
attachment budgets are checked in
\path{code/A_cubic_scalar_two.py} and its output.
\end{proof}
